\subsection{What the literatures leave open, and what this study assumes}
\label{sec:threat}

Three themes bear on this study, and each stops one step short of the same measurement.

\textbf{Why exposure of a field matters.} Identification does not need a name. Sweeney (2000) shows
that ZIP code, birth date and sex in combination identify most of the US population uniquely. El
Emam et al.\ (2011), reviewing
published re-identification attacks on health data, find that attacks succeed far less often against
data de-identified to a recognised standard than against data that was not. Between them these
establish that quasi-identifiers moving between components are a documented risk, which is why a
count of Safe Harbor fields crossing an agent boundary is worth taking. Neither supplies the
inverse: no mapping in either turns such a count into a probability of
re-identification, so a reduction in exposure remains a reduction in exposure.

\textbf{Where a control can be placed.} The unit each control chooses sets the floor on what can be
withheld. Progent (Shi et al., 2025) puts policy at the tool-call boundary, governing what an agent
may \emph{do}, not what it may see. Jacob et al.\ (2025) come closest, separating privilege by data type so an injected instruction
cannot travel with a value, but the unit is type, not field. Neupane et al.\ (2025) score
sanitisation and policy against a HIPAA least-privilege requirement without measuring inter-agent
handoffs. Together they establish that run-time confinement is buildable and auditable against a
regulatory standard. All three stop above the field, and none reports what the confinement it
imposes costs the task.

\textbf{What escapes when nothing is enforced.} This theme measures leakage without controlling it,
and its two members agree on where it concentrates. AgentLeak (El Yagoubi et al., 2026) reports 27.2 per
cent leakage in final outputs against 68.8 per cent in inter-agent messages, so output-only audits
miss most of it. MAGPIE (Juneja et al., 2025) finds up to 50.7 per cent leakage across two
hundred collaborative tasks, under explicit instruction not to disclose. Instruction alone is
therefore not a sufficient control, and enforcement belongs at a boundary. What it yields is a
channel and a rate. Both count what escapes without asking what the task loses when it is held
back.

The three give a risk worth measuring, a boundary to enforce at and a channel to measure on. Among
the sources reviewed we found none that does both: confine field by field on a clinical
data-quality task, then report the cost.

This study's threat model concerns \emph{mistakes}, not attackers: an agent holding patient data it
never needed through a bug or misconfiguration. Projection (Section~\ref{sec:intro}) rules out one
such failure, since an absent field cannot be read. Four survive.

\begin{itemize}[leftmargin=1.4em, itemsep=0.15em]
  \item \textbf{Over-granting.} A manifest may release more than the check needs, and projection
        faithfully delivers it. Enforcement is only as good as the declaration, and the ablation
        (Section~\ref{sec:ablation}) caught one such field.
  \item \textbf{Contents of a granted field.} A named field's contents may exceed the need, free text
        the plain case, since field granularity is not content granularity.
  \item \textbf{Onward leakage.} An agent may leak what it was legitimately given, to a log or a
        later tool call, which projection can only shrink.
  \item \textbf{Composition.} Two manifests, each minimal, may have projections that jointly
        re-identify a patient, a composition this study neither models nor bounds.
\end{itemize}

Every surviving failure is a \emph{consumer} failure, not a producer one. The producer side is the
larger omission: \texttt{project()} reads the whole record before building the reduced copy. The
tiered design gathers identifiers into two departments of its own, concentrating identifier access
into a single high-value target no adversary model here prices (Section~\ref{sec:limits}). Prompt
injection and side channels against the projection layer are out of scope. Projection bounds what an
injected instruction reaches, not whether it executes. A mistake that never crosses that bound falls
outside the metric.
